% ============================================================================
% Fichier cible : ARIEL_LOGMINER_MEMOIRE_FINAL/main.tex, chapitre étoilé « Bokuse »
% Lignes à remplacer : 517 à 626 (de « Mémoire oyo ezali kolakisa Ariel
%   \logminer{} » jusqu'à « mpo na analyse ya ba journaux hétérogènes. »).
%   Conserver les lignes 513-515 et le bloc « Maloba ya ntina » (628-636).
% Auteur : Agent 1 (intégrité factuelle), chantier 2026-11-05.
% Portée : seuls les chiffres et affirmations factuellement faux ou absents du
%   corps sont corrigés, pour aligner le Bokuse sur main.tex__resume_fr.tex.
%   Les paragraphes non concernés sont repris À L'IDENTIQUE. Les phrases
%   nouvelles (marquées [N] ci-dessous en commentaire) suivent le registre
%   existant (lingala + termes techniques français) mais N'ONT PAS été
%   vérifiées par un locuteur natif : relecture linguistique requise.
% ============================================================================

Mémoire oyo ezali kolakisa Ariel \logminer{}, architecture moko ya
multi-agents ya pete mpe oyo ekabolami na environnement ya laboratoire
mpo na kosala analyse ya ba journaux ya système mpe ya réseau oyo ezali
ya lolenge ndenge na ndenge. Agent moko ezali awa entité logicielle oyo
ezali na makoki, état, mémoire mpe politique locale na yango ; ekoki
kopesa proposition ya mosala, koboya yango, mpe na sima kosala yango na
nzela ya ba handlers oyo ekokani na mosala yango. Négociation kati na
ba agents etongami likolo ya protocole Contract Net oyo elimbolami
polele, mpe transport ekoki kozala local to kosalema na nzela ya Redis.
Autonomie mpe distribution ezali awa kolakisa bizaleli ya logiciel oyo
emekami na laboratoire. Intelligence générale mpe fonctionnement ya SOC
industriel ezali na libanda ya périmètre ya mosala oyo.

% [N] paragraphe modifié : retrait de l'écart-type, de l'IC et du F1 DDoS
% (absents du corps du texte) ; « F1 macro » -> « F1 moyen ».
Évaluation emonisi mingimingi influence ya protocole expérimental likolo
ya performances. Na CICIDS2017, Random Forest oyo emekami na séparation
aléatoire stratifiée ezwi F1 moyen ya $0{,}995142$ mpo na ba graines
mitano. Tango scénario moko na moko kati na ba scénarios mitano
elongolami na entraînement, F1 moyen ekiti kino $0{,}157744$.
Comparaison ya ba modèles légers mitano na holdout oyo emonisi ete
Logistic Regression ezwi F1 moyen ya malamu koleka, oyo ezali
$0{,}233670$, nzokande HistGradientBoosting ezwi PR-AUC macro ya
$0{,}569002$ kasi F1 na yango ekomi pene na zéro na seuil oyo eponamaki.
Ba résultats oyo ezali kolakisa ete kobongola modèle ekoki kobongola ba
métriques mosusu, kasi elongolaka te fragilité oyo emonanaka kati na ba
scénarios.

% [N] paragraphe modifié : BGL, unité corrigée (événements et non templates)
% et ajout de la coïncidence avec la règle de sévérité (corps, section 4.4).
Mpo na HDFS mpe BGL, protocole ekabolaka na ndenge ya makasi entraînement,
validation mpe test. Drain3 ebongisami kaka na entraînement mpe na sima
ekangami liboso ya kosalela ba partitions mibale oyo elandi. Na niveau
\texttt{block\_id}, Histogram ekomi na F1 ya $0{,}868373$ (IC 95~\% :
$[0{,}856771\,;\,0{,}879493]$) na ba blocs $115\,013$ ya test HDFS,
kati na yango $1\,679$ ezali positifs. Na BGL, F1 ezali $0{,}913698$, kasi
$89{,}8081\,\%$ ya ba événements ya test ezali na template oyo eyebanaki
te na train, mpe décision ya Histogram ekokani mobimba na règle
« sévérité absente de l'entraînement normal ». Na yango, score oyo emekaka
nouveauté ya sévérité, kasi modélisation ya dynamique ya ba journaux te ;
règle « template inconnu » ekomi kaka na $0{,}277885$.

% [N] paragraphe modifié : raccourci par le port de destination (section 4.6.4).
Na CSE-CIC-IDS2018, séparation kati na mokolo ya 15 mpe ya 16 février
epesi F1 moyen ya $0{,}999212$ na régression logistique contre
$0{,}401375$ na Random Forest. Audit emonisi nzokande ete
\texttt{Dst Port} yango moko ekomi na $0{,}999720$, score oyo règle fixe
likolo ya port de destination (na nse ya 1024) ezwi mpe, mpe ete transfert
ya 16 kino 15 février eleki te $0{,}234195$. Na yango, score principal
etongami likolo ya raccourci oyo ezali kaka ya mokolo ya test.

% [N] paragraphe modifié : remplissage du schéma, fixture Apache, absence de
% comparateur sans routage (tableau 7.1, H3).
Validation multiformat etangi, eparse mpe enormalise
$7\,001/7\,001$ unités, kati na yango fixture moko kaka ya Apache.
Résultat oyo ezali komeka couverture fonctionnelle ya ba formats oyo
emekamaki ; remplissage ya schéma commun ekeseni na source : ba champs
standard 7 na 7 mpo na syslog mpe BGL, kasi ata moko te mpo na ba flux
réseau tabulaires. Politique open-set oyo ecalibramaki
kaka na ba sources oyo eyebani, mpe na sima ekangami, eboyi $20/20$
fichiers inconnus mpe eboyi na libunga $3/46$ fichiers connus ; ba fichiers
connus 43 oyo endimami etambwisami malamu, mpe marge ya décision ezali te
probabilité calibrée. Comparaison na chaîne oyo ezali na routage te
esalemaki te ; yango wana gain prédictif ya routage spécialisé endimisami te.

% [N] paragraphe modifié : coût d'orchestration et panne appariée
% (sections 4.9 et 4.9.1) ; retrait de la liste « litomba architecturale ».
Benchmark contrôlé esangisi architectures minei, charges minei mpe
répétitions zomi, elingi koloba $160$ runs mpe $264\,000$ traitements
kozanga échec. Na ba microtâches ya protocole, monolithe eleki, mpe ba
architectures ya ba agents ebakisaka coût ya orchestration. Na panne
appariée oyo esalemi nsima ya persistance, ba architectures ya ba agents
ezongeli effet moko te, kasi monolithe mpe workers centralisés ezongeli
yango ; soki magasin d'idempotence moko wana ebakisami na yango, ba
doublons elongwe mpe.

% [N] paragraphe remplacé : la campagne multi-VM décrite est celle du corps
% (endurance d'une heure, section 4.11, tableau A.4) ; saturation (4.11.3).
Campagne d'endurance ya ngonga moko, na ba machines virtuelles mibale
(Debian mpe Ubuntu) oyo esangisami na Redis centralisé na hôte ya
laboratoire, esali $1\,800/1\,800$ tâches ; Debian ezwi 451 mpe Ubuntu
$1\,349$, engebene na ba profils de capacité oyo ebongisamaki. Arrêt
contrôlé nsima ya persistance ememi na réattribution moko mpe replay
kozanga effet dupliqué. Na saturation, capacité ya topologie oyo ezali
kati na 8 mpe 16 tâches/s, mpe machine ya mibale esali kaka
$1{,}23\,\%$ ya ba tâches oyo elongaki : ebakisi capacité te. Haute
disponibilité emekamaki te.

% [N] paragraphe modifié : ajout du corpus étendu (section 4.12).
Corrélateur, oyo emekamaki na corpus synthétique contrôlé, ezwi précision
ya $0{,}764706$, rappel ya $0{,}866667$ mpe F1 pairwise ya
$0{,}812500$. Na ba corpus générés ya minene, rappel na yango ekiti kino
$0{,}5973$ ; tango 30~\% ya ba incidents esangisi ba catégories
d'événements ebele, variante oyo etalaka kaka hôte mpe fenêtre temporelle
ezwi F1 ya malamu koleka ($0{,}7268$ contre $0{,}6496$).

Na campagne CNP multi-source, condition M esali mpenza ba
artefacts compatibles na $1400/1401$ unités, na fallback Apache moko
kaka mpe kozanga erreur. Latence moyenne na yango ezali
$0{,}105420$~s, contre $0{,}031027$~s mpo na condition heuristique H.
Lokola ba labels homogènes ezali te, ekoki kosalema te calcul ya
exactitude globale multi-source ; F1 Linux/auth ya $0{,}390244$ etikali
kaka résultat descriptif ya source oyo eponamaki.

Dépôt ezali mpe na procédure ya comparaison ya ba modèles mpe mode
dry-run. Ba tests ezali ko couvrir ba branches mosusu, kasi ezali ko
couvrir te cycle mobimba mpe oyo ezongelami mingi ya promotion mpe ya
rejet. Litomba globale ya procédure oyo emekamaki te.

% [N] paragraphe modifié : statut conforme au tableau 7.1 actuel
% (provisoire, modifiable par l'Agent 2).
Hypothèse moko te endimisami mobimba : misato (H1, H2 mpe H4) esimbami
kaka na ndambo, mpe mibale (H3 mpe H5) etikali NON TRANCHÉE, mpo
comparateur sans routage ezali te mpo na H3, mpe gain quantitatif ya
mémoire emonisami te mpo na H5. Biyano ya ba questions de recherche
etikali partielles ; question méthodologique (QR6) nde ezwi eyano oyo
ekoki koleka na kati ya périmètre ya laboratoire.
Contribution principale ya Ariel \logminer{} ezali bongo na intégration
architecturale mpe méthodologique ya chaîne locale, modulaire
mpe oyo ba expériences na yango ekoki kolandelama na ndenge ya polele
mpo na analyse ya ba journaux hétérogènes.
